\section{Source corrections and the scope of the conclusions}
\label{audit:sec:corrections}

The following audit concerns the pinned manuscript, rather than earlier
versions already corrected there.  It distinguishes a branch error, an
unsupported description of a special value, and an editorial clarification
about coefficient fields.  The accompanying patch changes only the first
two items; it does not alter the source files in this research package.

\subsection{A real branch correction for the arctanh integral}

In \path{chapters/03-algebraic.tex}, equation
\texttt{cleo:eq:lintanh} occurs in the discussion of real parameters of
absolute value less than one.  It writes
\[
 \int_0^{\pi/2}\operatorname{arctanh}(\lambda\sin x)\,dx
 =\theta\log\tan\frac\theta2
       +2\operatorname{Cl}_2(\theta)
       -\frac12\operatorname{Cl}_2(2\theta),
 \qquad \theta=\arcsin\lambda.
\]
The displayed logarithm is correct for positive $\lambda$.  For negative
$\lambda$, its argument is negative, and the principal complex logarithm
adds an imaginary term that cannot occur in the real integral.

\begin{proposition}[The real-parameter formula]
\label{audit:prop:arctanh}
For $-1\le\lambda\le1$, put $\theta=\arcsin\lambda$.  Then
\begin{equation}
 \int_0^{\pi/2}\operatorname{arctanh}(\lambda\sin x)\,dx
 =\theta\log\left|\tan\frac\theta2\right|
       +2\operatorname{Cl}_2(\theta)
       -\frac12\operatorname{Cl}_2(2\theta).
 \label{audit:eq:arctanh-corrected}
\end{equation}
At $\lambda=0$ the right side is its continuous value zero.  At
$\lambda=\pm1$ the integrals are convergent improper integrals.
\end{proposition}

\begin{proof}
Let the integral be $F(\lambda)$.  For $|\lambda|<1$, differentiation
under the integral is valid locally uniformly in $\lambda$.  The
substitution $t=\cos x$ gives
\[
 F'(\lambda)=\int_0^1
       \frac{dt}{1-\lambda^2+\lambda^2t^2}
 =\frac{\arcsin\lambda}{\lambda\sqrt{1-\lambda^2}},
\]
where the quotient has value one at zero.  Therefore
$dF(\sin\theta)/d\theta=\theta/\sin\theta$.

The classical Fourier definition of the Clausen function gives
$\operatorname{Cl}_2'(u)=-\log|2\sin(u/2)|$ away from multiples
of $2\pi$.  On either interval $(-\pi/2,0)$ or $(0,\pi/2)$,
differentiating the right side of \eqref{audit:eq:arctanh-corrected}
gives $\theta/\sin\theta$: the three logarithmic terms cancel by
$\sin\theta=2\sin(\theta/2)\cos(\theta/2)$.
Both sides tend to zero at the origin, proving the identity on
$(-1,1)$.  Finally $\operatorname{arctanh}(\sin x)$ has only an
integrable logarithmic singularity at $x=\pi/2$; dominated convergence
and oddness give the endpoint cases.
\end{proof}

The failure at a specific admissible negative parameter is exact.  At
$\lambda=-1/2$, $\theta=-\pi/6$, the printed principal-logarithm
right side differs from the integral by
\[
 i\pi\theta=-\frac{i\pi^2}{6}.
\]
The real integral is approximately
$-0.5317299165586139653837501727942413753$.
Thus the correction is a branch issue, not a numerical uncertainty.
The smallest repair is to insert the absolute value in the logarithm
and specify the continuous value at zero.  Restricting the printed
identity to $0<\lambda<1$ would also make its original logarithm valid.

\subsection{Remove a remaining unsupported non-elementarity label}

The same chapter, at line 194 of the pinned source, calls
\[
 \int_0^{\pi/2}\arctan^2(\sin x)\,dx
       =\pi\chi_2(3-2\sqrt2)
\]
the ``simplest non-elementary diagonal.''  The exact evaluation is
retained.  The manuscript gives neither a specified comparison algebra
nor a proof that this constant lies outside one.  Its later discussion
correctly distinguishes exact formulas from unproved minimality and
non-elementarity.  The proposed patch replaces the unsupported adjective
by an explicit statement that no elementary reduction or nonreduction
theorem for this value is established there.  This is a correction of
the claim's status; it is not a proof of a new reduction.

\subsection{Coefficient fields depend on the coordinates}

The paragraph following \texttt{integral:res:parity} in
\path{chapters/07-integration.tex} says that the reduction coefficients
live in $\mathbb Q(\sqrt q)$.  That description fits the immediately
displayed real quadratic character modulo five.  It should not be read
as a field assertion for arbitrary rational grids or arbitrary
characters.  Chapter~2 already correctly uses the cyclotomic DFT field
$\mathbb Q(\zeta_q)$; passing to multiplicative characters can also
introduce character values and Gauss sums.

For example, the primitive character modulo five specified by
$\chi(2)=i$ takes the values $1,i,-i,-1$ on $1,2,3,4$.
The character-coordinate coefficients therefore include $i$, which is
not in the real field $\mathbb Q(\sqrt5)$.  Even a real sine-coordinate
presentation need not stay in that quadratic field:
\[
 \sin^2\frac{2\pi}{5}=\frac{5+\sqrt5}{8},\qquad
 N_{\mathbb Q(\sqrt5)/\mathbb Q}
      \left(\sin^2\frac{2\pi}{5}\right)=\frac5{16}.
\]
If $\sin(2\pi/5)$ belonged to $\mathbb Q(\sqrt5)$, the last
quantity would be the square of a rational norm, which it is not.
The displayed quadratic-character identity itself requires no repair.
The appropriate clarification is to state the coefficient field for
the chosen additive, character, or real-paired coordinates.

\subsection{Prior in-repository continuations must be credited}

The pinned repository already contains substantial periodic Stieltjes
calculus outside the manuscript chapters, under
\path{docs/reports/stieltjes-correlation-calculus/continuations/}
relative to \path{Analysis/Polylogarithms/} \cite{RepoConvolutionAlgebra,RepoConvolutionCalculus,RepoShiftedJets}.
In particular, \path{convolution-algebra/article.tex} proves the
one-generator Bell algebra in \texttt{thm:Bell}, the Gamma-quotient
multiplication law in \texttt{thm:product}, the polygamma contact
correction, and the reflected spectral kernel in
\texttt{thm:reflected}.  The file
\path{convolution-calculus/Stieltjes_Convolution_Calculus.tex}
already records the symmetric first-Stieltjes order-zero
polylogarithm formula in \texttt{eq:Li0identity}.
The file \path{shifted-hurwitz-jets/sections.tex} gives the complete
normalized primitive tower in \texttt{shj:Uformula} and proves its
endpoint regularity.

The untwisted periodic results identified above therefore consolidate and
independently verify those untwisted foundations. The nonintegral-twist
section separately answers the explicit open target in the archived
convolution-algebra report.  The appearance of a
question in an earlier intake report does not establish its current open
status when a later in-repository continuation has already answered it.
The shifted harmonic Laurent and nonintegral-twist results are assessed
separately, with their own source and literature qualifications, rather
than deriving a novelty claim from the untwisted periodic construction.

An explicitly open target in the same
\path{convolution-algebra/article.tex} is the subsection
``Twists, characters, and alternate endpoint normalizations.''
It asks for the nonintegral frequency-shift analogue and an all-index
half-twist table with its point-supported terms.
Section~\ref{tw:sec:main} answers this specified target: the convolution
unit is $\delta_0$, the scalar Gamma coefficients remain ordinary zeta
values, the basis consists of Lerch jets, and the covariant contact
terms are given explicitly by Theorem~\ref{tw:thm:contacts}.
The nonsingular alternating convolution and the modified Stieltjes
functions themselves have the antecedents cited there.  The resolved
repository question concerns their singular extension and normalization.

\subsection{What this continuation does not prove}

At the pinned snapshot, the mixed sums $S_2$ and $S_4$ already have
exact proofs.  The weight-seven relation
\texttt{cycloquot:conj:S6} and the newer weight-nine relation
\texttt{s8new:conj:S8} remain conjectural.  A different frozen
weight-nine vector was rigorously rejected in
\texttt{research:prop:S8-rejected}; its rejection must not be confused
with the status of the newer candidate.  No result of the present
continuation promotes either surviving numerical candidate to a theorem.

The distributional identities proved above fix their extensions through
meromorphic continuation, or equivalently the stated cutoff finite parts.
Agreement with an ordinary function away from the endpoint and zero mean
alone do not characterize every higher derivative extension: derivatives
of the endpoint delta distribution have zero mean.  The explicit endpoint
terms in those theorems are part of the statements and must be preserved
when the formulas are integrated or differentiated.
